\section{Pre-registration record and deviations}\label{app:prereg}
\setcounter{table}{0}\setcounter{figure}{0}\setcounter{equation}{0}

% generated by scripts/make_cmame_material.py from the records; do not edit
\begingroup\small
\begin{longtable}{@{}>{\raggedright\arraybackslash}p{0.14\linewidth}>{\raggedright\arraybackslash}p{0.42\linewidth}>{\raggedright\arraybackslash}p{0.21\linewidth}>{\raggedright\arraybackslash}p{0.11\linewidth}@{}}
\caption{Every pre-registered criterion of the runs reported here, with its outcome.}\label{tab:criteria}\\
\toprule
Run & Criterion & Measured & Outcome \\
\midrule
\endfirsthead
\caption[]{Every pre-registered criterion of the runs reported here, with its outcome. (continued)}\\
\toprule
Run & Criterion & Measured & Outcome \\
\midrule
\endhead
\bottomrule
\endfoot
2D, 16 July 2026 & G1$'$ (a), energy-term sub-experiment (separately trained networks): supervised + energy term at least 3 times better than the zero field in displacement, and better than the 3 naive baselines fitted on 1,024 labelled instances, at every budget & 4.5--6.1 times; naive baselines beaten at every budget & met \\
 & G1$'$ (b), Tables 1 and 2, 64 labels: the energy term lowers the relative energy gap by $\geq$ 50\% and the von Mises error by $\geq$ 25\% & 78.7\%, 65.4\% & met \\
 & G1$'$ (c), Tables 1 and 2, 64 labels: label-free, then supervised beats supervised by $\geq$ 10\% in relative energy gap or $\geq$ 5\% in displacement & 92.6\%, 2.0\% & met \\
 & Gate G1$'$ = (a) and (b) and (c) &  & GO \\
 & K1, energy-term sub-experiment: the better of the fixed and the gradient-scaled energy term lowers the relative energy gap by $<$ 25\% at every budget & 98.0\% / 94.4\% / 90.6\% / 74.0\% & not triggered \\
 & K2, Table 1: label-free displacement error $>$ 30\% above supervised, 1,024 labels & +3.6\% & not triggered \\
 & C1-advantage, Table 2: label-free relative energy-gap advantage over supervised $<$ 40\% at every budget & 98.7\% / 97.0\% / 93.5\% / 80.4\% & not triggered \\
 & K4: a latent-space regulariser raises the standardised effective rank of the latent vectors by $\leq$ 1.5 times & 1.70 times & not triggered \\
 & K5: conjugate gradients from the prediction of a label-free network (64 validation instances) save $<$ 20\% of the iterations from zero & $-$0.9\% & triggered \\
 & K6: a cross-resolution invariance term shrinks the transfer gap by $<$ 10\% at coarsening 2.5 & $-$23.7\% & triggered \\
 & E6 (alignment): within-geometry rank correlation of latent and solution distances $<$ 0.3 & 0.730 & not triggered \\
 & C5: any pre-registered check of the conditioning bound, the modewise identity or the conjugate-gradient bound violated & none & not triggered \\
\midrule
2D, 31 July 2026 & K3: a self-supervised joint-embedding term added to label-free pretraining changes the displacement error and the relative energy gap of the transformer, fine-tuned with 16, 64 and 256 labels, by less than 3\% at every budget (if triggered, the term is dropped) & change with the term, 16 / 64 / 256 labels: displacement $-$13.5\% / $-$0.5\% / +2.6\%; relative energy gap $-$48.4\% / +92.5\% / $-$16.4\% (positive: more accurate) & not triggered \\
\midrule
3D, 21 September 2026 & Initial run: gate G2 and kill conditions as registered; the label-free objective was defective (deviation D14) & (a), (b), (c) not met; KP1, KP2, KP4 triggered & NO-GO \\
3D, 22 September 2026 & Pilot of the amendment: label-free displacement error after 20 epochs $<$ 0.90 & 0.309 & met \\
3D, 28 September 2026 & G2 (a), budgets $\geq$ 64 (amendment), Table 2: supervised + energy term at least 3 times better than the zero field in displacement, and better than both naive baselines & 3.9 / 12.0 / 21.3 times; naive beaten & met \\
 & G2 (a), every budget (the registered form, reported for reference) & 2.71 times at 16 labels & not met \\
 & G2 (b), Tables 2 and 3: label-free displacement error within +10\% of supervised at 1,024 labels, and relative energy-gap advantage $\geq$ 40\% at every budget & $-$36.3\%; 99.95\% / 99.85\% / 98.86\% / 99.97\% & met \\
 & G2 (c), Table 4: finer mesh, zero-shot displacement error $\leq$ 1.25 times in-band, and better than both naive baselines & 8.62 times; best naive 0.401 against 0.258 & not met \\
 & Gate G2 = (a) and ((b) or (c)), amended &  & GO \\
 & Gate G2, registered form (reference) &  & NO-GO \\
 & KP1: supervised (1,024 labels) better than label-free by $>$ 10\% in displacement & $-$36.3\% & not triggered \\
 & KP2: relative energy-gap advantage $<$ 40\% at any budget & minimum 98.9\% & not triggered \\
 & KP3: the energy term lowers the relative energy gap by $<$ 25\% at every budget & +96.6\% / +85.7\% / $-$20.7\% / +99.6\% & not triggered \\
 & KP4: finer mesh, zero-shot displacement error $>$ 1.5 times in-band, or a naive baseline better & 8.62 times & triggered \\
 & KP5: any pre-registered check of the conditioning bound, the modewise identity or the conjugate-gradient bound violated & none & not triggered \\
 & KP6 (alignment): within-geometry rank correlation of latent and solution distances $<$ 0.3 & 0.827 & not triggered \\
\midrule
2D, 29 September 2026 & Latent regularisation, K1: relative change of the displacement error or of the relative energy gap beyond the resolution & within the resolution & not triggered \\
 & Latent regularisation, K2: the separation does not rise in any seed & rose in every seed & not triggered \\
 & Latent regularisation, GO: separation raised by $\geq$ 0.02 in every seed & +0.009 / +0.014 / +0.002 & NO-GO \\
\midrule
3D, 30 September 2026 & Token bottleneck, M = 512, K1: in-band relative energy gap or finer-mesh displacement error worse than the transformer's by more than the resolution & +280.9\% (resolution 287.2\%), +79.1\% (resolution 41.9\%) & triggered \\
 & Token bottleneck, M = 512, K2: finer-mesh training step longer than 2 s & 0.049 s & not triggered \\
 & Token bottleneck, M = 512: outcome &  & dropped \\
 & Token bottleneck, M = 1,024, K1: in-band relative energy gap or finer-mesh displacement error worse than the transformer's by more than the resolution & +183.3\% (resolution 70.5\%), +122.0\% (resolution 33.3\%) & triggered \\
 & Token bottleneck, M = 1,024, K2: finer-mesh training step longer than 2 s & 0.047 s & not triggered \\
 & Token bottleneck, M = 1,024: outcome &  & dropped \\
\end{longtable}
\endgroup
\vspace{-\smallskipamount}
\noindent\begin{minipage}{\linewidth}\footnotesize\raggedright
Labels are those of the pre-registration documents, each of which numbers its own criteria. Gates (G1$'$, G2) decide whether the study proceeds to its next stage. Kill conditions withdraw a claim, or abandon a component, when triggered: in two dimensions K1--K6, C1-advantage (the claimed advantage in relative energy gap), C5 (the pre-registered checks of \cref{sec:theory:checks}) and E6 (latent alignment); in three dimensions KP1--KP6, with KP5 the counterpart of C5. The latent-regularisation and token-bottleneck runs number their criteria afresh, with the meanings given here. A criterion reads the trained networks of the tables it names; the energy-term sub-experiment of the 2D run of 16 July 2026 trained its own supervised networks with and without the energy term, with the same configuration. K4, E6 and KP6 concern latent representations, which this paper does not use. Standardised effective rank: the participation ratio of the eigenvalues of the covariance of the latent vectors, each dimension standardised. The 3D amendment, the deviations and the joint-embedding term are described below, the latent regularisation in \ref{app:e1} and the token bottleneck in \cref{sec:cost:training}.
\end{minipage}
\par\medskip


\Cref{tab:criteria} lists every criterion fixed in advance for the runs reported in this
paper, with the value measured and the outcome as recorded by the code at the time. The
two-dimensional pre-registration also retired, before its run, an earlier criterion (a
displacement improvement of at least 10\% from the energy term with 64 labels) that an
earlier run had already answered in the negative; the run of 16 July 2026 measured
\numGOneRetired{} with the fixed-weight form of the energy term, and the criterion played no
part in its decisions. The following departures from the plans bear on the results. They
are recorded in the repository's list of deviations and its provenance note, except the
remark at the end of the paragraph on the energy term, which was noticed while this paper
was written.

\paragraph{Deviation D14}
The initial three-dimensional run (21 September 2026) completed its registered protocol and
read NO-GO, with conditions (a), (b) and (c) of its gate unmet and kill conditions KP1,
KP2 and KP4 triggered. After the verdict, its label-free networks were found to have a
displacement error of \numDfourteenDisp{} and a relative energy gap of
\numDfourteenGap{}, identical to five decimals across seeds and consistent with predicting
the solution scaled by $\alpha=\numDfourteenAlpha{}$ (displacement error $1-\alpha$,
relative energy gap $(1-\alpha)^2$). The cause was in the code: inference multiplied the
decoded field by $F_{\max}$, the label-free objective scored it without that factor, and on
this corpus $F_{\max}\approx\alpha$. The fix gives the network one decoding path, used by
the objective and by inference, with a test that fails on the registered code. The
supervised networks, which are trained and evaluated through the inference path, were not
affected. The verdict of the initial run stands as recorded. A pre-registered amendment then
repeated everything that depended on the label-free objective: the label-free networks,
their fine-tuning on the finer mesh and the latent-alignment criterion read from them. It
reused the 30 trained supervised networks and the 6 networks trained from scratch on the
fine mesh, each checked against its recorded SHA-256 hash, and it required a short pilot
run first: one seed, 20 epochs, which had to bring the displacement error below
\numPilotLimit{} and reached \numPilotDisp{}. A first version of the amendment was
registered and withdrawn on the same day, before any run, because the audit of its
execution path had not been completed; the second version was registered after the audit.
The amended run (28 September 2026) was executed once, without interruption.

\paragraph{The sanity condition of the three-dimensional gate}
Condition (a) of the three-dimensional gate required the supervised transformer with the
energy term to be at least three times more accurate in displacement than the zero field,
and better than both naive baselines, at every budget. In the initial run it failed at 16
labels (\numGTwoSanitySixteen{} times), independently of D14. The amendment, written with
this result known, assessed condition (a) only from the decision budget of 64 labels, where
it was met by factors of \numGTwoSanitySixtyFour{} to \numGTwoSanityMax{}. Its stated
rationale is that the factor of three had been calibrated on the two-dimensional family,
and that on the three-dimensional family the supervised network itself sits below it with
16 labels and above it from 64 labels on. In two dimensions the condition had read the
networks of the energy-term sub-experiment, which met it at every budget; the networks of
\cref{tab:labeleff} would not have met it with 16 labels (\numGOneSanityTableSixteen{}
times), because of the scatter described in \cref{sec:results2d:labels}. The amendment
declared the change as made after the fact and required the original form to be computed
and reported beside it; both readings are in \cref{tab:criteria}. The change affects only
the gate's verdict, not any measurement reported in this paper.

\paragraph{Deviation D16}
The text of the three-dimensional pre-registration states that the label-free network
trains on the 1,744 in-band instances outside the validation set. The registered
configuration, which is what ran, trained it on the first 1,024 of them, the same instances
on which the supervised networks were trained with labels. Every three-dimensional run
executed this way. The discrepancy was found after the amended run and is recorded with the
other deviations. It works against the label-free network, which saw fewer instances than
the text promised, and it leaves the cleaner comparison: the same instances, with and
without their solutions.

\paragraph{The energy term with 16 labels}
In the networks of \cref{tab:2d,tab:labeleff}, the two-dimensional supervised transformer
with the energy term used the gradient-scaled form at every budget, and with 16 labels its
displacement errors scattered across seeds (\cref{sec:results2d:labels}); the energy-term
sub-experiment of the same run trained both the gradient-scaled and the fixed-weight form.
An exploratory run (30 July 2026), with six seeds in two separate experiments, gave
displacement errors from \numDiagScaledMin{} to \numDiagScaledMax{} for the gradient-scaled
form, against \numDiagFixedMin{} to \numDiagFixedMax{} in all \numDiagFixedSeeds{} seeds for
a fixed weight of one on the normalised energy. The three-dimensional pre-registration
therefore fixed that weight at 16 labels and the gradient-scaled form from 64 on. The
description of the protocol stored in the three-dimensional report still names the
gradient-scaled form for every budget; the registered code applied the pre-registered
policy.

\paragraph{The joint-embedding term (criterion K3)}
Criterion K3 of the two-dimensional pre-registration concerned a separate run
(31 July 2026, \ref{app:repro}). That run pretrained the transformer without labels for 200
epochs, once on the energy alone and once on the energy together with a self-supervised
joint-embedding term (\ref{app:e1}), which predicts the latent vectors of masked regions of
the mesh from those of the rest, with a regulariser on the mean latent vectors; it then
fine-tuned both with 16, 64 and 256 labels for 200 epochs, with three seeds. K3 would have dropped the term, and with it the version of the
planned paper that used it, had the term changed neither the displacement error nor the
relative energy gap by more than \numKthreeBand{} at any budget. K3 was not triggered, and
the registered rule therefore kept the term (\cref{tab:criteria}). The criterion did not
state a direction. Of the \numKthreeOutside{} changes outside the band, \numKthreeWorse{}
made the fine-tuned network less accurate; the remaining one, a relative energy gap
\numKthreeGain{} lower with 64 labels, compares with a baseline whose seeds scattered
widely (relative energy gaps of \numKthreeBaseSeeds{}, against \numKthreeBaseRef{} for the
same configuration in the run of 16 July 2026). This paper reports the energy objective
without the term. That scope was chosen after the result was known, and it departs from
the version of the paper that the rule selected.

\paragraph{Provenance of the two-dimensional runs}
The two-dimensional runs were executed outside a version-controlled checkout, so their
reports record no commit. The code checked their configuration hashes against the
pre-registration documents at start-up, and the code version was established afterwards
from the runs' console output and by comparing the source trees of the released packages.
The pre-registration documents of these runs were committed to the repository on 3 August
2026, after the runs; the repository's provenance note, compiled on 1 August 2026, records
the hashes of the documents that the runs checked, and those of the reports and of the
archives that contained them. The three reports used here are committed with their
hashes (\ref{app:repro}).

\paragraph{Audit}
An independent script recomputed the three-dimensional verdict from the report and agreed
on \numAuditItems{} items. Its first pass reported one discrepancy, which was in the audit
script itself: it had combined the gate's conditions more strictly than the pre-registered
rule. Both outputs are archived.
